\documentclass[a0,portrait,ruledsections]{sciposter}
%
\usepackage[utf8]{inputenc}
\usepackage[T1]{fontenc}
\usepackage{babel}
\usepackage{lmodern}
\renewcommand*\familydefault{\sfdefault}
\usepackage{fourier}
\usepackage[protrusion=true,expansion=true]{microtype}	% Better typography

\usepackage{multicol}
\usepackage{graphicx}
\usepackage{svg}
\DeclareGraphicsExtensions{.pdf,.png,.jpg}
\usepackage{blindtext}
\usepackage[absolute]{textpos}
%\usepackage{biblatex}
%% Nice units
\usepackage{siunitx}
\usepackage[autostyle=true,%
            german=guillemets,%
            english=british]{csquotes}


%% To include Tikz from octave/matlab
\usepackage{pgf}
\usepackage{tikz}
\usepackage{pgfplots}
\pgfplotsset{compat=newest}
\pgfplotsset{plot coordinates/math parser=false}


% === Color definitions
\definecolor{BoxCol}{cmyk}{0.95,0.50,0.10,0.25}	% HSB Blau
\definecolor{SectionCol}{cmyk}{0.95,0.50,0.10,0.25} % HSB Blau

\renewcommand{\titlesize}{\Huge\textbf}
\renewcommand{\authorsize}{\normalsize}
\renewcommand{\instsize}{\large}
%
\conference{\textbf{\copyright ~Hochschule Bremen, School of Electrical Engineering and Computer Sciences, MScEE, 2025}} 
%
\setmargins[2cm]
%
\rightlogo[0.75]{./figures/HSB_Vertikal_RGB}
\leftlogo[0.65]{./figures/qr_code}
\title{Assessing the feasibility of developing and taping-out an analog block for LED driver control using IHP-Open-PDK
  and IIC-OSIC-TOOLS}
\author{P. Toyni and H. Pangavhane}
%   
\institute{Prof. Dr.-Ing. M. Meiners, Concept Engineering SOCs and Design}
\email{mirco.meiners@hs-bremen.de}

%%%%%%%%%%%%%%%%%%%%%%%%%
\begin{document}

\maketitle

%\vspace{1cm}
% 
\begin{multicols}{3}
%
\begin{abstract}
Open-source EDA has reached a point where its practical viability for real mixed-signal IC fabrication can be examined on modern semiconductor processes. This thesis evaluates that potential through the design and tape-out of an analog block for LED-driver control using the IHP OpenPDK SG13G2 a 130 nm SiGe BiCMOS PDK and the implementation using the open-source IIC-OSIC-TOOLS developed by Prof. Pretl (JKU Linz).

A key contribution of this work is practical verification: two tape-outs were accepted and fabricated in the IHP Open MPW run of July 2025. To strengthen cross-disciplinary understanding, the authors executed two complete tape-outs, one for a design of NMOS foldedcascode structure and other for PMOS foldedcascode structure, switching roles one serving as design engineer while the other served as layout engineer in the first run, and vice versa in the second. This approach revealed workflow bottlenecks, toolchain strengths, and realistic guidance for analog teams adopting open-source EDA.

The results demonstrate that a fully open-source toolchain, combined with an industry-grade SiGe BiCMOS open PDK, can reliably produce manufacturable analog ICs with predictable post-layout performance. This work serves as a practical reference for engineers, researchers, and SMEs evaluating open-source flows for mixed-signal IC development, and provides evidence that democratized, open analog IC design is not only feasible but tape-out proven.
\end{abstract}
%
%=======================================================================================================


% -----------------------------------------------------------------------------

\section{Motivation \& Scope}\label{sec:ubersicht}

\textbf{Industry-Grade Analog Tape-Out Using Open-Source EDA}

Analog IC development in industry is defined by strict timelines, verified
design flows, and manufacturable silicon. This work follows the same
constraints. Using the IHP SG13G2 130\,nm SiGe BiCMOS process, a complete
analog design flow was executed—from schematic and layout to verification
and fabrication—entirely with open-source EDA tools.

The objective was not to propose a new methodology, but to prove feasibility
under realistic industrial conditions. Two independent analog designs were
completed, verified, and submitted within a standard MPW schedule, resulting
in successful tape-out and fabrication. The results demonstrate that
open-source EDA can satisfy industrial requirements for analog IC design,
verification, and delivery.

% -----------------------------------------------------------------------------


\section{Design and Fabrication Context}

The analog blocks were designed and fabricated using the IHP SG13G2
130\,nm SiGe BiCMOS technology, selected for its industrial relevance,
robust device models, and availability as an open PDK. The complete
design flow was executed using the open-source IIC-OSIC-TOOLS toolchain,
covering schematic capture, simulation, layout, parasitic extraction,
and verification.

Both designs followed a standard industry-style MPW schedule and were
submitted to the IHP Open MPW run of July~2025. All verification steps,
including DRC, LVS, and post-layout simulation, were completed prior to
submission. The flow and timeline were aligned with industrial analog IC
development practices.

% -----------------------------------------------------------------------------

\section{Analog Block Architecture}

The target circuit is a folded-cascode operational transconductance
amplifier (OTA), selected as a representative analog building block for
LED driver control applications. The topology offers high intrinsic gain,
good output swing, and robust biasing characteristics, making it suitable
for integration in mixed-signal power management systems.

Two functionally equivalent implementations were realized: an NMOS
folded-cascode variant and a PMOS folded-cascode variant. This enabled a
direct comparison of device-level trade-offs and layout complexity within
the same technology node and verification flow. Both designs were
dimensioned for stable operation under post-layout parasitics and
industrial operating conditions.

\begin{figure}
\centering
\includesvg[width=0.85\linewidth]{./figures/_fig_foldedcascode_diagram.svg}
\caption{Folded-cascode OTA topology implemented in NMOS and PMOS variants.}
\label{fig:foldedcascode}
\end{figure}

% \vspace{-1cm}

% -----------------------------------------------------------------------------

\section{Dual Tape-Out Strategy and Role Separation}

To evaluate the robustness of the open-source design flow under realistic
conditions, two independent tape-outs were executed within the same MPW
cycle. The first tape-out implemented an NMOS folded-cascode OTA, while
the second implemented a PMOS folded-cascode OTA, both using identical
verification and submission procedures.

To reflect industrial team structures, the authors deliberately separated
responsibilities between design and layout. In the first tape-out, one
author served as circuit designer while the other was responsible for
layout and verification; these roles were reversed for the second tape-out.
This approach enabled an unbiased assessment of workflow handover,
toolchain usability, and verification effort.

The dual tape-out strategy exposed practical bottlenecks and validated the
repeatability of the open-source flow, demonstrating that independent
analog designs can be delivered within an industry-style MPW timeline.

% -----------------------------------------------------------------------------

\section{Layout, Verification, and Post-Layout Validation}

Both OTA variants were fully laid out following analog layout best
practices, including device matching, symmetry, and substrate isolation.
Guard rings and well-defined supply routing were applied to ensure robust
operation under mixed-signal conditions.

Design rule checking (DRC) and layout-versus-schematic (LVS) verification
were completed successfully for both designs. Parasitic extraction (PEX)
was performed, and the extracted netlists were used for post-layout
simulation. The post-layout results showed predictable behavior relative
to pre-layout simulations, confirming that the open-source flow supports
reliable parasitic-aware analog design.

\begin{figure}
\centering
\includegraphics[width=0.85\linewidth]{./figures/_fig_otaled}
\caption{Folded-cascode OTA topology implemented in NMOS and PMOS variants.}
\label{fig:foldedcascode}
\end{figure}

% -----------------------------------------------------------------------------

\section{Results and Key Takeaways}

Both analog designs were successfully verified and fabricated through the
IHP Open MPW program, demonstrating that a fully open-source EDA flow can
support industry-grade analog IC development. The results confirm that
complex analog blocks can be designed, verified, and delivered within a
standard MPW timeline without reliance on proprietary tools.

This work validates the repeatability and practical usability of open-source
EDA for analog and mixed-signal design and provides a proven reference
flow for academic groups, SMEs, and early-stage industrial teams targeting
real silicon fabrication.









%=======================================================================================================

%\bibliography{../_bibliography/references.bib}

\end{multicols}

\end{document}
